% Options for packages loaded elsewhere
\PassOptionsToPackage{unicode}{hyperref}
\PassOptionsToPackage{hyphens}{url}
%
\documentclass[
]{article}
\usepackage{amsmath,amssymb}
\usepackage{iftex}
\ifPDFTeX
  \usepackage[T1]{fontenc}
  \usepackage[utf8]{inputenc}
  \usepackage{textcomp} % provide euro and other symbols
\else % if luatex or xetex
  \usepackage{unicode-math} % this also loads fontspec
  \defaultfontfeatures{Scale=MatchLowercase}
  \defaultfontfeatures[\rmfamily]{Ligatures=TeX,Scale=1}
\fi
\usepackage{lmodern}
\ifPDFTeX\else
  % xetex/luatex font selection
\fi
% Use upquote if available, for straight quotes in verbatim environments
\IfFileExists{upquote.sty}{\usepackage{upquote}}{}
\IfFileExists{microtype.sty}{% use microtype if available
  \usepackage[]{microtype}
  \UseMicrotypeSet[protrusion]{basicmath} % disable protrusion for tt fonts
}{}
\makeatletter
\@ifundefined{KOMAClassName}{% if non-KOMA class
  \IfFileExists{parskip.sty}{%
    \usepackage{parskip}
  }{% else
    \setlength{\parindent}{0pt}
    \setlength{\parskip}{6pt plus 2pt minus 1pt}}
}{% if KOMA class
  \KOMAoptions{parskip=half}}
\makeatother
\usepackage{xcolor}
\setlength{\emergencystretch}{3em} % prevent overfull lines
\providecommand{\tightlist}{%
  \setlength{\itemsep}{0pt}\setlength{\parskip}{0pt}}
\setcounter{secnumdepth}{-\maxdimen} % remove section numbering
\ifLuaTeX
  \usepackage{selnolig}  % disable illegal ligatures
\fi
\usepackage{bookmark}
\IfFileExists{xurl.sty}{\usepackage{xurl}}{} % add URL line breaks if available
\urlstyle{same}
\hypersetup{
  hidelinks,
  pdfcreator={LaTeX via pandoc}}

\author{}
\date{}

\begin{document}

\section{Intrinsic finite-energy Monge--Ampere products in nef
classes}\label{intrinsic-finite-energy-mongeampere-products-in-nef-classes}

This paper states and proves the exact nef operator and
variational-support theorem independently accepted in Xean review124.
The field is algebraically closed, trivially valued, of characteristic
zero and arbitrary cardinality; the original variety is smooth integral
and projective of dimension \(n\geq1\). The datum is any actual nef
rational line bundle \(L\), with no semiampleness or finite-generation
assumption. All potentials and measures remain on the entire original
\(X^{\mathrm{an}}\).

The reference is the canonical zero metric. The raw mass is
\(V=c_1(L)^n\). On the entire genuine numerical cone define

\[E_\theta(u)=\inf_{k\geq1}E_\theta(\max(u,-k)),\qquad
\mathcal E^1(\theta)=\{u\in\operatorname{PSH}(\theta):E_\theta(u)>-\infty\}.\]

The theorem constructs a target-independent positive Radon top operator
and common-class mixed products of mass \(V\) on this whole intrinsic
domain, including endpoints outside the auxiliary ample finite-energy
vector space. It proves reference-floor total-variation convergence,
integrability of every intrinsic test, weighted arbitrary decreasing-net
convergence against all original continuous tests, both signs of the
original continuous-direction outer derivative, weak upper
semicontinuity and compact energy-cost sublevels. At \(V>0\), the
intrinsic finite-dual condition is equivalent to integration of every
genuine finite-energy test, every constructed output satisfies it, and
its variational objective is coercive.

At \(V=0\), the positive top and common-class mixed operators, energies
and weighted products vanish. No probability normalization, inverse
equation or uniqueness is claimed in this range. This paper proves no
all-target attainment or uniqueness theorem at positive volume either.
Thus its reviewed scope is a supporting theorem for original P16.5; it
does not close that original problem, and supplies no general nonnef
big-class extension.

The only supplied equation input is the exact actual-ample G0 theorem,
used solely at positive actual ample perturbations. The published
model-pairing, envelope and numerical inputs are named in each proof
part and remain in their stated domains. They are imported mathematical
foundations, not claims of recursively proving every cited theorem. The
four proof parts below preserve the frozen mathematical bodies
independently reconstructed in review124. Their earlier draft-status
sentences and chronological remaining-obligation paragraphs are retained
as provenance; the present theorem and Part V specify the accepted final
scope.

\subsection{0. The ample input and the finite-continuous closure
replacement}\label{the-ample-input-and-the-finite-continuous-closure-replacement}

The supplied ample theorem concerns an actual ample rational bundle on
the retained original characteristic-zero projective variety. It
constructs its genuine finite-energy metric operator, raw energy and
continuous-obstacle envelopes, supplies compactness and all-point
divisorial determination, and proves both continuous-direction outer
energy derivatives. It gives no assertion for a nonample class. Its
exact full mathematical source is the previously accepted P16.1 paper;
the independently frozen review packet includes that entire body. Within
the present proof every invocation has an actual ample class, such as
\(\theta+\epsilon h\) with \(\epsilon>0\).

For clarity, the closure step avoids an arbitrary extended-valued
flattening lemma. Genuine nef membership equals the intersection of the
genuine cones for \(\theta+\epsilon h\). The forward section embeddings
preserve the scalar, since generating canonical section norms have
maximum zero. For the reverse direction choose finite continuous ample
FS majorants of a potential in every positive perturbed cone, allowing
positive constant shifts. They form a single directed decreasing family:
for any finite collection of old continuous majorants, their finite
continuous minimum is an upper bound; with a positive strict tolerance,
eventual approximation below that minimum holds on an open cover, and
compactness provides one common later approximant. Refine the positive
perturbation at the same step. Infimizing tolerances and majorants gives
the desired original genuine all-point limit, including its poles. This
uses finite continuous old minima, never a general varying-class
flattening principle for arbitrary discontinuous functions. Closed
intersections in one normalized ample compact cone provide the genuine
normalized nef compactness and the envelope passage used below.

\subsection{I. Bounded genuine metric products and continuous
obstacles}\label{i.-bounded-genuine-metric-products-and-continuous-obstacles}

This is a new unreviewed own deduction for original P16.5. It is not
supplied to either active frozen native119 or native121. Original
closure and completion counts do not change. The deduction advances the
general nef case without semiampleness, while preserving the distinction
between bounded potentials and the larger unbounded nef finite-energy
class.

\subsubsection{1. Actual scope and the corrected source
boundary}\label{actual-scope-and-the-corrected-source-boundary}

Let k be algebraically closed, trivially valued and of characteristic
zero, with arbitrary cardinality. Let X be smooth, integral and
projective of dimension n\textgreater=1. Retain an actual nef rational
line bundle L, its numerical class theta, and its canonical trivial
reference zero. Fix an actual ample rational bundle H with class h. All
potentials below belong to the genuine numerical cone of BJ
Definition4.1; no quotient cone or inverse-lift operator is defined.

The primary source is Boucksom--Jonsson, arXiv:1801.08229v4, header
March2022, 143pages. Theorem5.20 itself states the envelope property for
\textbf{every nef theta}, not just ample theta, under the retained
smoothness and characteristic assumptions. Its printed proof reduces to
ample rational classes through Lemma5.9; that reduction does not narrow
the theorem's conclusion. Complete printed statement/proof units checked
here are local literal LF2980--2985 and3027--3062, and fresh native-PDF
lines4269--4274 and4321--4363. Lemma5.9's complete proof is already
checked at LF2801--2842. Published recursive inputs in the
multiplier-ideal argument are retained, not independently recursively
certified.

Two distinctions matter. Lemma5.17's equivalence with continuity of
continuous-obstacle envelopes is formulated for ample classes.
Theorem5.20 does not by that citation alone prove continuity of every
nef continuous-obstacle envelope. Remark7.16 further warns that the
floor-defined E1(theta) can strictly contain the ambient finite-energy
vector space intersected with PSH(theta). Theorems7.14/7.18 and
Proposition7.19 therefore cannot simply be applied to every member of
E1(theta).

The exact whole accepted ample G0 proof is
\path{/workspace/research/root-checks/public-P16.1-round-92-final-reader/argument.md},
SHA256
\path{6cd26d798b63a388bd9a5e6dcd0583a7cd9c7417648efe390f4d4d00d28646b9}.
Only its actual ample projective domain is a grant. Its raw energy,
genuine metric operator, energy concavity, continuous envelopes, and
both continuous-direction outer derivatives are used at theta+epsilon h.
No nonample equation or unbounded nonample continuity is granted.

\subsubsection{2. Bounded genuine potentials fit the actual published
product
domain}\label{bounded-genuine-potentials-fit-the-actual-published-product-domain}

Write B(theta)=PSH(theta) intersect bounded functions. This is a
statement about the entire genuine numerical bounded cone, not a newly
defined smaller test family. Since theta is nef, zero and every constant
are genuine theta-psh by BJ4.3. If u belongs to B(theta), then u belongs
to PSH(theta+h), since zero is h-psh and numerical addition is BJ4.7. It
is bounded, so its ample energy is finite by monotonicity between its
constant bounds. Thus u belongs to the actual ambient vector space
E\textasciitilde1 of Definition7.12.

Theorems7.14 and7.18 now apply with their stated hypotheses. Define the
raw bounded energy and mixed measures by their actual energy-pairing
formulas:

\[E_\theta(u)=\frac{(\theta,u)^{n+1}-(\theta,0)^{n+1}}{n+1},\qquad
M_\theta(u_1,\ldots,u_n)=\bigwedge_{i=1}^n(\theta+dd^cu_i).\]

The total mass is theta\^{}n, and M is positive for genuine bounded
theta-psh inputs, by Theorem7.18 and Proposition7.19(i). The energy
shifts by c theta\^{}n under u -\textgreater{} u+c.~At a finite PL/model
stage, Proposition7.19(ii) gives the actual sum of component Dirac
measures with coefficients

\[\operatorname{ord}_F(\mathcal X_0)\prod_i(\theta_{\mathcal X}+D_i)|_F.\]

These are actual model metric intersection products with the given
reference, component multiplicities and bundle class. On continuous
potentials the published CLD compatibility of Remark7.23 applies.
Individual current factors are not asserted to have been constructed on
all bounded inputs merely from the wedge notation.

One can also recover the same bounded theta-products as limits of actual
positive ample operators before taking any limit. For
epsilon\textgreater0 rational, each u\_i remains genuinely
(theta+epsilon h)-psh. The actual G0 ample operator on L+epsilon H is
defined, positive, and has raw mass (theta+epsilon
h)\^{}n.~Multilinearity of the published mixed measures makes its
pairing with each continuous test a finite polynomial in epsilon, with
the constant term M\_theta(u\_1,\ldots,u\_n). Thus the positive ample
products converge weakly to the displayed genuine theta-product.
Uniqueness in Theorem7.18 makes this independent of H. The limit is
identified with an already constructed original metric product; it is
not the definition of an arbitrary inverse lift.

For a decreasing net u\_i in B(theta) with bounded genuine limit u, any
tail has common upper and lower bounds. All terms lie in E1(theta+h),
and Proposition7.19(iv) yields the original bounded mixed weak
continuity. Theorem7.14 gives bounded-energy continuity for these nets.
Integrability and the energy-pairing identity for an ambient
E\textasciitilde1 test are Proposition7.19(v). This proves no weighted
theorem for a test outside that stated ambient domain.

\subsubsection{3. Genuine envelopes at the nef
boundary}\label{genuine-envelopes-at-the-nef-boundary}

For continuous f on the original compact X\^{}an, let P\_theta f be the
pointwise supremum of genuine theta-psh functions bounded by f.~Constant
competitors show that it lies between min f and max f.~The nef envelope
property makes the upper regularization theta-psh; continuity of f makes
that regularization \textless=f, so it is itself a competitor. Hence
P\_theta f equals that regularization, is genuinely theta-psh, and is
bounded. This argument does not assert its continuity.

Put theta\_epsilon=theta+epsilon h and q\_epsilon=P\_theta\_epsilon
f.~The actual ample theory makes q\_epsilon continuous. As epsilon
decreases to zero, q\_epsilon decreases pointwise. The common constant
bounds make the limit generically finite, and BJ(PSH2) makes its limit
theta-psh. That limit is \textless=f.~Conversely every theta-psh
competitor is theta\_epsilon-psh, so P\_theta f\textless=q\_epsilon.
Therefore

\[q_\epsilon\downarrow q_0=P_\theta f.\]

All q\_epsilon, for epsilon\textless=1, lie in the one actual ample cone
PSH(theta+h) and have the same two constant bounds. Theorems7.14/7.18
and Proposition7.19 therefore apply to this decreasing family within
their proper domain. Polynomial expansion, coefficient by coefficient,
proves

\[E_{\theta_\epsilon}(q_\epsilon)\longrightarrow E_\theta(q_0),\qquad
M_{\theta_\epsilon}(q_\epsilon)\longrightarrow M_\theta(q_0).\]

Some individual theta-coefficients at positive epsilon may be signed;
positivity of those coefficients is neither used nor asserted. Their
continuous-test pairings converge by Proposition7.19(iv) in the common
ample finite-energy cone. The full operators on the left and the final
operator on the right are positive. Every finite polynomial coefficient
has a finite limit, so the positive powers of epsilon vanish without an
unproved uniform estimate over arbitrary E1(theta).

\subsubsection{4. Uniform continuity needed for
variation}\label{uniform-continuity-needed-for-variation}

Pointwise order and constants give the original envelope estimate

\[\|P_\theta f-P_\theta g\|_\infty\leq\|f-g\|_\infty.\]

For uniformly converging bounded genuine theta-psh potentials u\_j
-\textgreater{} u, their theta mixed measures and energy converge. Here
is a reduction to the published decreasing-net theorem that avoids any
unproved nonample strong topology. From any sequence converging in sup
norm, pass to a subsequence with errors \textless=4\^{}\{-j\}; adding
3\emph{4\^{}\{-j\} makes that subsequence decreasing pointwise to u.
Constants do not change curvature products. Energy changes only by
3}4\^{}\{-j\} theta\^{}n.~The decreasing convergence in the common ample
cone from Section2 therefore gives the desired convergence on that
subsequence. Every subsequence admits such a further subsequence, so the
entire sequence converges. For each fixed continuous measure test this
proves continuity on the sup-norm metric space; consequently it also
proves convergence for arbitrary sup-norm converging nets. It does not
apply a false monotone convergence theorem to arbitrary Borel nets.

\subsubsection{5. Both derivatives of continuous-obstacle
energy}\label{both-derivatives-of-continuous-obstacle-energy}

Define F\_theta(f)=E\_theta(P\_theta f) for f in C0(X\^{}an). It is
finite, concave and translation equivariant with raw coefficient
theta\^{}n.~Concavity follows from the genuine convex cone, envelope
order and bounded energy concavity; alternatively it follows from the
pointwise limit of the ample concave functionals in Section3.

We first give the needed ample supporting inequalities rather than
importing an unproved assertion about arbitrary obstacles. For an ample
class omega and continuous f, q=P\_omega f is continuous genuine. If
0\textless=s\textless=1, monotonicity gives P\_omega(q+s(f-q))=q. The
actual ample G0 outer derivative at q in the continuous direction f-q
therefore yields orthogonality:

\[\int(f-q)M_\omega(q)=0.\]

Energy concavity along genuine psh segments gives
\[
 E_\omega(r)\le E_\omega(q)+\int(r-q)M_\omega(q)
\]
for \(r=P_\omega g\). Since \(r\le g\) and \(M_\omega(q)\) is positive, the displayed
orthogonality proves the supporting inequality

\[F_\omega(g)-F_\omega(f)\leq\int(g-f)M_\omega(P_\omega f).\]

Use this inequality once at f and once at f+tg. For t\textgreater0 it
bounds the quotient between integral g M\_omega(P\_omega(f+tg)) and
integral g M\_omega(P\_omega f). Apply it with omega=theta\_epsilon and
pass epsilon down to zero using the actual energy and operator limits
from Section3. Thus

\[\int gM_\theta(P_\theta(f+tg))\leq
\frac{F_\theta(f+tg)-F_\theta(f)}t\leq
\int gM_\theta(P_\theta f).\]

The envelope Lipschitz estimate and Section4 make the left integral
converge to the right one as t down to zero. For t\textless0 the same
two supporting inequalities reverse under division; the two bounds have
the same limit as t increases to zero. Therefore both derivatives exist
and agree:

\[\left.\frac{d}{dt}\right|_{t=0}F_\theta(f+tg)=\int gM_\theta(P_\theta f),\qquad f,g\in C^0(X^{\rm an}).\]

In particular this proves both outer signs at every \textbf{continuous
genuine} theta-psh base potential u, by taking f=u. It proves the
continuous-obstacle formula for every nef actual L, without
semiampleness and including zero raw top mass. It does not infer a
derivative at an arbitrary unbounded member of the floor-defined
E1(theta), or even at an arbitrary bounded discontinuous base, from this
continuous-obstacle argument. There is no contradiction with
proposal131: that proposal uses a different lower-degree mixed energy of
positive mass, whereas the present top-degree energy has zero mass for
its product L and is constant.

\subsubsection{6. Remaining original
obligations}\label{remaining-original-obligations}

This supplies a genuine bounded top-degree operator, nef envelope
membership, and both derivatives for continuous obstacles beyond
semiample contractions. It does not supply the complete unbounded
E1(theta) operator with all original weighted/floor limits, intrinsic
target-space characterization, or variational attainment and uniqueness
for every finite-energy measure. The nef raw top mass is theta\^{}n;
proving all required section-volume and reference identifications in the
intended full general class remains part of the original problem.
General nonnef big references and general nonbig nonzero classes are
also still unresolved here.

The explicit published source domains are the important constraint.
Remark7.16 prohibits silently identifying the floor-defined E1(theta)
with E\textasciitilde1 intersect PSH(theta). A bounded approximation
theorem cannot be promoted to that larger class without an energy
estimate controlling the full limiting operator. All new conclusions
here require independent reconstruction. Original P16.5 remains open.

\subsubsection{7. Source/application
ledger}\label{sourceapplication-ledger}

\begin{itemize}
\tightlist
\item
  BJv4 LF2190--2237: genuine numerical definition, joint decreasing
  class closure, constants for nef theta. LF2330--2368: addition and
  uniform closure proof. Used for Sections2--4.
\item
  BJv4 LF2801--2842: Lemma5.9 and Theorem5.11 complete
  statements/proofs. LF2980--2985 and3027--3062: complete Theorem5.20
  statement and proof. Fresh native-PDF selected units are retained as
  source audit only. Used for nef envelope membership, not unproved nef
  continuity.
\item
  BJv4 LF2909--2935: complete Lemma5.17 statement/proof, whose ample
  domain is retained. Used only for theta\_epsilon continuous envelopes.
\item
  BJv4 LF3890--3955: actual ambient E\textasciitilde1 definition,
  complete Theorem7.14 and Lemma7.15 proofs, and Remark7.16.
  LF3980--4057: complete Theorem7.18 and Proposition7.19
  statements/proofs. Used only on verified bounded inputs in a common
  ample finite-energy cone.
\item
  Exact whole ample G0 body named in Section1: actual raw operator,
  energy concavity and both outer derivatives within its accepted ample
  domain. Every use at positive epsilon retains that domain. No original
  P16.5 or other unproved nonample premise is a grant.
\item
  Primary URL: https://arxiv.org/pdf/1801.08229v4 . Whole-book and
  recursive-foundation verification are not claimed. No finite numerical
  check certifies the deduction.
\end{itemize}

\subsection{II. The entire intrinsic floor operator and integrable
weights}\label{ii.-the-entire-intrinsic-floor-operator-and-integrable-weights}

This is a new unreviewed own deduction, not a grant or an original
completion. Neither active frozen native119 nor native121, nor prepared
frozen review122, receives this later material. The construction
addresses the larger floor-defined nef E1 class without assuming that
its members belong to an auxiliary ample finite-energy space. Arbitrary
decreasing-family continuity, a full variational inverse and general
original P16.5 remain open.

\subsubsection{1. Actual data and verified bounded input
domain}\label{actual-data-and-verified-bounded-input-domain}

Retain an algebraically closed trivially valued field k of
characteristic zero and arbitrary cardinality, a smooth integral
projective n-fold X with n\textgreater=1, an actual nef rational bundle
L with class theta, and its canonical reference zero. Put
V=theta\^{}n\textgreater=0. Use the genuine numerical PSH(theta) of BJv4
Definition4.1. Its constants and finite maxima are genuine. No
semiampleness is assumed.

For bounded u in PSH(theta), membership in PSH(theta+h) for an actual
ample H follows by adding the genuine zero h-potential. Boundedness
makes u finite-energy in that actual ample class, so u is in the
published ambient E\textasciitilde1. Theorems7.14/7.18 and
Proposition7.19 therefore construct its \textbf{actual} raw metric
measure

\[\nu(u)=(\theta+dd^cu)^n\]

and raw energy E\_theta(u)=(theta,u)\^{}\{n+1\}/(n+1), with zero
reference height. The measure is positive, of mass V, and its finite
model coefficients include ord\_F(X\_0) times the actual intersection
degree on F. Multilinearity identifies it with the epsilon-down-to-zero
limit of actual positive ample measures for L+epsilon H, independently
of H. This identification uses actual metric operators before any limit
and defines no arbitrary inverse measure. These bounded facts are
reconstructed in proposal132 and follow from the complete named
published input units just specified.

Normalize u in PSH(theta) by sup u=0. Its value at the generic trivial
valuation is then zero, as follows already by taking limits of the
defining FS weights. Let u\_j=max(u,-j), and define

\[E_\theta(u)=\inf_{j\geq1}E_\theta(u_j),\qquad
\mathcal E^1_\theta=\{u\in\operatorname{PSH}(\theta):E_\theta(u)>-\infty\}.\]

This is the \textbf{entire actual floor-defined cone}, including
potentials outside E\textasciitilde1. Finite energy is the original
theta-energy condition, with raw normalization. For u in that cone set
C\_u=-E\_theta(u)\textless infinity.

\subsubsection{2. Borel locality at the nef
boundary}\label{borel-locality-at-the-nef-boundary}

We need locality on the strict floor sets, not merely on ordinary open
subsets. For bounded theta-psh a,b, both are finite-energy theta+epsilon
h-psh. The complete actual ample Theorem7.40 yields equality of the two
actual ample metric measures of max(a,b) and a on the Borel set
\{a\textgreater b\}. Multiply by the raw positive ample volume to remove
probability normalization.

Each of these raw measures is a finite polynomial in epsilon as a signed
Radon measure, by the published product construction. Restriction to the
fixed Borel set is a linear operation on these finite signed measures.
Thus their restricted difference is a polynomial identically zero for
every positive rational epsilon. Evaluating against any Borel set, or
solving the finite Vandermonde system at distinct rational epsilon
values, makes every measure coefficient zero. Its constant coefficient
is the actual bounded theta-product. Hence

\[\mathbf1_{\{a>b\}}\nu(\max(a,b))=
\mathbf1_{\{a>b\}}\nu(a).\]

There is no passage of a merely weak limit through restriction to a
discontinuous indicator. This exact coefficient argument retains the
general Borel strict-wins statement of Theorem7.40, whose complete
printed proof was read at literal LF4531--4533 and4565--4586, with
Lemma7.43 LF4552--4564. The source's printed proof uses its
model-component locality and then the actual weighted finite-energy
passage. Its apparent reference to Proposition7.19(iii) at the end is a
source cross-reference issue; the substantive weighted passage is
Proposition7.21(ii), not the scaling assertion in (iii).

\subsubsection{3. Construct the full operator from actual
floors}\label{construct-the-full-operator-from-actual-floors}

Let A\_j=\{u\textgreater-j\}, nu\_j=nu(u\_j), and
sigma\_j=1\_\{A\_j\}nu\_j. For m\textgreater=j, max(u\_m,-j)=u\_j and
\{u\_m\textgreater-j\}=A\_j. Section2 gives
nu\_m\textbar A\_j=nu\_j\textbar A\_j. Since A\_j increases, the actual
sigma\_j form an increasing family of positive Radon measures. Borel
restrictions of finite regular Radon measures on the compact original
space are again regular Radon measures.

The bounded energy identity is the actual mixed formula

\[E_\theta(u_j)=\frac1{n+1}\sum_{r=0}^n
\int u_j\,(\theta+dd^cu_j)^r\wedge\theta^{n-r}.\]

Every mixed measure is positive. Since u\_j\textless=0, this implies

\[\int(-u_j)\nu_j\leq(n+1)(-E_\theta(u_j))\leq(n+1)C_u.\]

On A\_j\^{}c, u\_j=-j, so the mass of rho\_j=1\_\{A\_j\^{}c\}nu\_j is at
most (n+1)C\_u/j. Consequently sigma\_j has mass tending to V. Its
increasing setwise limit is a positive Radon measure sigma of mass V:
the difference sigma-sigma\_j has mass tending to zero, giving
regularity by finite-mass tail approximation. Moreover

\[\|\nu_j-\sigma\|_{\rm TV}\leq
\|\rho_j\|+\|\sigma-\sigma_j\|
\leq\frac{2(n+1)C_u}{j}.\]

Define M\_theta(u)=sigma. This is a measure obtained from the original
actual positive metric floor products and their strict-wins
restrictions. It has the required raw top mass V, gives no mass to
\{u=-infinity\}, and has total-variation convergence of the actual full
floors. For bounded u the construction agrees exactly with the published
bounded operator, since eventually A\_j is all of X\^{}an. All auxiliary
ample choices have already disappeared at each bounded stage.

The same construction with arbitrary real floor levels is increasing and
cofinal with the integers. Changing u by a constant merely shifts the
real levels and leaves curvature products unchanged. Thus the operator
and finite-energy membership are intrinsic under the stated constant
gauge. If V=0, all bounded positive products vanish and all bounded
energies vanish by their constant bounds. The resulting operator is zero
on the entire genuine cone; this does not supply uniqueness modulo
constants or a positive lower-degree replacement.

\subsubsection{4. Every original E1(theta) potential is
integrable}\label{every-original-e1theta-potential-is-integrable}

For bounded normalized theta-psh a,b, monotonicity of the published nef
pairing makes every (theta,a)\textsuperscript{r(theta,b)}\{n+1-r\}
nonpositive. Its all-zero reference height is zero. The identity

\[\int b\,\nu(a)=(\theta,b)(\theta,a)^n-(\theta,0)(\theta,a)^n\]

therefore gives -integral b nu(a) \textless=
-(theta,b)(theta,a)\^{}n.~Expanding (2theta,a+b)\^{}\{n+1\}, all its
terms are nonpositive and the displayed mixed term occurs n+1 times.
Concavity of bounded energy gives

\[\int(-b)\nu(a)\leq-E_{2\theta}(a+b)
=-2^{n+1}E_\theta((a+b)/2)
\leq2^n(-E_\theta(a)-E_\theta(b)).\]

The concavity here follows from the actual ample concavity and
polynomial passage on these bounded inputs; it is not applied to an
unverified nonample ambient domain.

Apply the estimate to a=u\_j and b=v\_l, where u,v are any normalized
original E1(theta) potentials. It bounds integral(-v\_l)nu\_j by
2\^{}n(C\_u+C\_v), uniformly in j,l. First take j to infinity, using
total-variation convergence with the bounded Borel test v\_l. Then let l
increase by ordinary countable monotone convergence. Thus

\[\int(-v)M_\theta(u)\leq2^n(C_u+C_v)<\infty.\]

In particular the constructed measure gives no mass to the pole set of
\textbf{any} original E1(theta) test. This is an actual integrability
statement on the larger intrinsic domain, not an assumed embedding into
E\textasciitilde1.

\subsubsection{5. Strong weighted limits along the actual floor
sequence}\label{strong-weighted-limits-along-the-actual-floor-sequence}

Suppose V\textgreater0; the zero-mass case is immediate. At bounded
theta-psh inputs, define the normalized I and J functionals by the
actual raw energies divided by V and the actual metric measures divided
by V. The fundamental energy identity, the bounds
I/(n+1)\textless=J\textless=nI/(n+1), and Corollary7.31's Holder
estimate transfer from theta+epsilon h by their finite polynomial and
pairing formulas. The constants remain bounded: by the source's explicit
convention at LF473 they depend only on n, while the positive volumes
tend to V\textgreater0. Every input in this passage is bounded and thus
lies in the actual ample finite-energy domain. No inequality on the
larger unbounded cone is assumed.

Let v\_l\textless=v\_k be two floors of a normalized v in E1(theta),
with l\textgreater=k, and set
delta\_kl=(E\_theta(v\_k)-E\_theta(v\_l))/V. The mixed energy difference
formula makes

\[\int(v_k-v_l)\nu(v_l)/V\leq(n+1)\delta_{kl},\qquad
I_\theta(v_k,v_l)\leq n(n+1)\delta_{kl}.\]

All normalized J-values of v\_k,v\_l,u\_j are bounded by
B=max(C\_v,C\_u)/V. Corollary7.31, applied with reference u\_j, gives

\[\int(v_k-v_l)\nu_j/V\leq\delta_{kl}
+C_n\big(n(n+1)\delta_{kl}\big)^{\alpha_n}B^{1-\alpha_n},
\qquad\alpha_n=2^{-n}.\]

Indeed J\_\{u\_j\}(v\_k)-J\_\{u\_j\}(v\_l) equals -delta\_kl plus that
integral. Put delta\_k=(E\_theta(v\_k)-E\_theta(v))/V. It tends to zero;
delta\_kl\textless=delta\_k. Letting l increase and using integrability
yields a bound integral(v\_k-v)nu\_j \textless=V r\_k, uniformly in j,
with r\_k -\textgreater0 given by the right side. The same bound holds
for M\_theta(u), by first passing j to infinity with bounded v\_k-v\_l,
and then l to infinity.

For fixed k, v\_k nu\_j tends to v\_k M\_theta(u) in total variation.
The two tail bounds consequently give

\[v\,\nu_j\longrightarrow v\,M_\theta(u)
\quad\hbox{in total variation}.\]

The same argument gives v\_j nu\_j -\textgreater{} v M\_theta(u) in
total variation. Thus all fixed original E1(theta) weights, as well as
their simultaneous original floors, have the claimed strong
\textbf{floor-sequence} limit. This statement retains the original
operator and weights before limits; it is stronger than merely checking
numerical mass.

\subsubsection{6. Precise remaining
boundary}\label{precise-remaining-boundary}

This construction supplies the actual entire floor-defined nef E1
operator, full raw mass, gauge and auxiliary-ample independence,
integrability of every intrinsic original E1 test, and the stated
weighted countable-floor limits. It does not establish arbitrary
decreasing nets of E1(theta) operators against every original continuous
test, arbitrary model approximations, or the required full variational
equation and uniqueness. The source's Holder estimates on theta-psh test
differences need not separate all continuous functions when theta is
nonample; replacing the original test space by their span would be a new
observable-carrier restriction. Those gaps remain explicit.

No general nonnef big reference theorem or positive lower-degree nonbig
carrier is constructed. A zero top operator on a nonconstant cone cannot
provide uniqueness modulo constants. Original P16.5 remains open, and
the new deduction requires an independent review before any scoped
acceptance count changes.

\subsubsection{7. Source/application
ledger}\label{sourceapplication-ledger-1}

\begin{itemize}
\tightlist
\item
  BJv4: genuine numerical definitions and joint closure; complete
  Theorems7.14/7.18, Lemma7.15 and Proposition7.19 at the selected units
  recorded for proposal132. These are applied only to the verified
  bounded inputs.
\item
  Actual ample G0 body and its raw concavity/metric-product
  compatibility: the exact whole accepted file named in proposal132,
  with verified SHA256
  6cd26d798b63a388bd9a5e6dcd0583a7cd9c7417648efe390f4d4d00d28646b9. Its
  nonample inverse extension is not a grant.
\item
  BJv4 literal LF4074--4125: complete Proposition7.21 and Lemma7.22
  statements/proofs, used to identify the substantive weighted passage
  in the printed locality proof. They are not applied directly to an
  arbitrary intrinsic E1(theta) potential.
\item
  BJv4 LF4175--4208: complete Proposition7.27 statement/proof.
  LF4247--4277 and4280--4294: complete Lemma7.30 and Corollary7.31
  statements/proofs. LF473: constants convention depends only on
  dimension. These are applied at positive ample epsilon on bounded
  inputs, before the proved nef coefficient limit.
\item
  BJv4 LF4531--4533,4552--4586: complete Theorem7.40 and Lemma7.43
  statements/proofs. The Borel polynomial-coefficient argument in
  Section2 is the own extension to bounded nef inputs.
\item
  Primary URL https://arxiv.org/pdf/1801.08229v4 . The actual floor
  construction and its limits are own proofs. Whole-book and
  recursive-foundation verification are not claimed.
\end{itemize}

\subsection{III. Arbitrary mixed weighted nets and both outer
signs}\label{iii.-arbitrary-mixed-weighted-nets-and-both-outer-signs}

This is an unreviewed own continuation. It combines the actual
bounded-domain reconstruction132, the intrinsic floor construction133
and the independently written native121 Section5. It is not an
independent acceptance of those inputs. Original P16.5 remains open.
Neither active native119 nor native122 receives this later material.

\subsubsection{1. Domain and actual starting
constructions}\label{domain-and-actual-starting-constructions}

Retain an algebraically closed trivially valued characteristic-zero
field of arbitrary cardinality, smooth integral projective X of positive
dimension n, actual nef rational L, theta=c1(L), canonical reference0
and raw V=theta\^{}n.~No semiampleness or bigness is assumed. All
potentials belong to the original genuine numerical single-net cone
PSH(theta). The entire intrinsic E1(theta) is defined by the actual
bounded floors and their raw energy, as in133, including endpoints
outside the published auxiliary-ample signed E\textasciitilde1.

The bounded original products are constructed from actual metric
products with their model central multiplicities before any limit. Their
nef positivity and auxiliary independence use verified bounded
membership in E\textasciitilde1 and actual ample perturbations. Borel
winning-set locality uses exact polynomial coefficients, not weak
convergence through indicators. For normalized u in intrinsic E1(theta),
the original floor operator M(u) is constructed by increasing
winning-set restrictions and satisfies

\[\begin{aligned}
 M(u)(X^{an})&=V,\qquad C_u=-E_\theta(u),\\
 \|M(\max(u,-k))-M(u)\|_{TV}&\le 2(n+1)C_u/k.
\end{aligned}\]

It integrates every intrinsic E1(theta) test, with the estimate proved
in133. Its fixed-test and simultaneous literal-floor weighted limits are
in total variation. No nonample equation or arbitrary-model limit is a
starting grant.

\subsubsection{2. Energy and top continuity for the original arbitrary
decreasing
net}\label{energy-and-top-continuity-for-the-original-arbitrary-decreasing-net}

Let u\_i decrease pointwise on all of X\^{}an to a genuine generically
finite u. All suprema are bounded on a tail, because they are values at
the generic trivial valuation. At fixed k, max(u\_i,-k) decreases to
max(u,-k) in the verified bounded domain. Bounded energy continuity
gives convergence of the floor energies. Monotonicity yields

\[E_\theta(u)\le\inf_i E_\theta(u_i)
\le E_\theta(\max(u,-k))\quad\hbox{for every }k.\]

Taking k to infinity proves E\_theta(u\_i) decreases to E\_theta(u),
including the extended value -infinity. Simultaneous floors then give
concavity on the entire intrinsic E1 cone; constants shift raw energy by
Vc. The cone is upward finite-energy closed.

If the endpoint is intrinsic E1, choose a common constant B so that
u\_i-B\textless=0 on a tail. Their negative raw energies are bounded by
C=-E\_theta(u-B). The actual floor total-variation error is consequently
at most 2(n+1)C/k uniformly over that original tail and its endpoint.
For each original continuous f, first choose k to make this error small,
and then choose a late original index using bounded-floor weak
continuity. Thus

\[M(u_i)\longrightarrow M(u)\]

weakly against every f in C(X\^{}an). This retains the original
arbitrary-cardinality net. It uses no replacement of the continuous test
space by theta-psh differences and no countable subsequence of the
original net. A sequence is used only for the independent cofinal
reference floors.

This is the useful new observation in native121 Section5.4, now combined
with133: a uniform total-variation floor estimate resolves the earlier
concern about insufficient theta-test separation.

\subsubsection{3. Transfer the energy inequalities to intrinsic
E1}\label{transfer-the-energy-inequalities-to-intrinsic-e1}

Assume V\textgreater0 and divide the actual energies and measures by V
only in the auxiliary normalized I,J notation. Constants do not affect I
or relative J. Every bounded theta-psh input belongs to actual ample
E1(theta+epsilon h); published Proposition7.27 and
Lemma7.30/Corollary7.31 pass to theta by the actual finite-polynomial
coefficient limit. Their dimension-only constants stay bounded as the
positive ample volumes approach V.

For any finite tuple of intrinsic E1(theta) potentials, take their
simultaneous literal floors. The energy terms converge by definition.
Every weighted integral in the I,J identities and Holder bounds
converges by133, applied separately to each original fixed potential and
its actual floors. Hence all those identities and inequalities extend to
this larger intrinsic domain, without placing its endpoints in
E\textasciitilde1. In particular, for a\textgreater=b in E1(theta), with
delta=(E\_theta(a)-E\_theta(b))/V,

\[I(a,b)\le n(n+1)\delta.\]

One way to see the bound before passage is the positive mixed
energy-difference identity: integral(a-b)M(b)/V\textless= (n+1)delta,
whence J\_b(a)\textless=n delta and the fundamental I,J comparison gives
the displayed bound. This reasoning needs order but does not require
equal suprema.

If a\_i decreases to a in E1(theta), Section2 gives
delta\_i-\textgreater0 and therefore I(a\_i,a)-\textgreater0. Normalized
J-values are bounded on a tail because sup a\_i and E\_theta(a\_i)
converge. The transferred Lemma7.30 thus controls integrals of any fixed
intrinsic theta-psh test difference against M(a\_i)-M(a), as well as the
transferred Corollary7.31's relative-J differences.

\subsubsection{4. Fixed and decreasing original E1 weights: full
weighted weak
convergence}\label{fixed-and-decreasing-original-e1-weights-full-weighted-weak-convergence}

Let u\_i decrease to u in E1(theta), mu\_i=M(u\_i), mu=M(u), and let v
be a fixed intrinsic E1(theta) weight. Subtract a fixed constant so
v\textless=0; write v\_k=max(v,-k). The transferred Lemma7.30 applied to
v\_k-0 and u\_i-u gives

\[\int v_k\,d\mu_i\longrightarrow\int v_k\,d\mu\]

for every fixed k. Section2 gives mu\_i-\textgreater mu weakly on the
full original continuous test space. Since -v\_k is a bounded
nonnegative lsc function, the actual primary Lemma7.22, with a constant
net of weights, turns those two facts into weak convergence of
(-v\_k)mu\_i to (-v\_k)mu. This step handles a bounded discontinuous
v\_k; weak convergence alone would not.

The same Corollary7.31 argument as in133 supplies uniform weight tails.
All normalized J-values of u\_i on the tail are bounded by a common B,
as are those of v and its floors. Set

\[\delta_k=(E_\theta(v_k)-E_\theta(v))/V\downarrow0.\]

For l\textgreater=k, apply the bounded or transferred relative-J
inequality to v\_k,v\_l with reference u\_i, and use the ordered I
bound. Passing l to infinity by ordinary countable monotone convergence
gives

\[\int(v_k-v)\,d\mu_i
\le V\left(\delta_k+C_n[n(n+1)\delta_k]^{2^{-n}}B^{1-2^{-n}}\right)=:Vr_k,\]

uniformly on the tail, with r\_k-\textgreater0. The same inequality
holds for mu. For a continuous test f, the error replacing v by v\_k is
at most \textbar\textbar f\textbar\textbar{} Vr\_k at either end. First
choose k and then the original index. Hence v mu\_i converges weakly to
v mu against every original continuous f.

Now also let v\_i decrease to v in intrinsic E1(theta), indexed by the
same arbitrary directed set (a common refinement handles independently
indexed nets). Their J-values are bounded and their raw energy
difference tends to zero. The identical relative-J estimate applied to
v\_i,v, with reference u\_i, gives

\[\int(v_i-v)\,d\mu_i\longrightarrow0.\]

The integrand is nonnegative; thus the total-variation norm of
(v\_i-v)mu\_i tends to zero. Combining this with the fixed-v conclusion
proves

\[v_i M(u_i)\longrightarrow vM(u)\]

weakly as signed Radon measures on the original carrier. Constants are
restored by unweighted weak continuity. Only ordinary countable monotone
convergence was used for fixed-endpoint reference floors. No
arbitrary-Borel-net monotone convergence is substituted for the actual
weighted argument.

\subsubsection{5. Entire mixed products in one intrinsic nef
class}\label{entire-mixed-products-in-one-intrinsic-nef-class}

For u\_1,\ldots,u\_n in E1(theta), their every nonempty subset average
a\_I=\textbar I\textbar\^{}\{-1\}sum\_\{i in I\}u\_i belongs to
E1(theta), by genuine convexity and raw energy concavity. Define the
mixed limit by polarization of the already constructed top operators:

\[M(u_1,\ldots,u_n)=\frac1{n!}
\sum_{\varnothing\ne I\subset\{1,\ldots,n\}}
(-1)^{n-|I|}|I|^nM(a_I).\]

This formula identifies an actual product limit; it does not choose a
measure in response to a target. Indeed for bounded literal floors
u\_\{r,k\}, it is exactly the finite multilinear identity for the actual
original bounded products. Each subset average of those floors decreases
to a\_I with intrinsic E1 endpoint. Section2 therefore shows that the
actual positive mixed floor measures converge to the displayed
expression. The limit is positive, has raw mass V, and is independent of
common cofinal floor levels by the same argument.

We also identify its usual increasing winning-set construction. For
bounded mixed entries, mixed locality follows from bounded top locality
applied to a+sum(lambda\_r c\_r) and b+sum(lambda\_r c\_r) on the fixed
Borel set \{a\textgreater b\}. The class is (1+sum lambda\_r)theta.
Positive rational lambda values and the finite Vandermonde coefficient
argument recover each mixed coefficient. Apply this one entry at a time
on W\_k=\{all u\_r\textgreater-k\}: later actual mixed floors have the
same restriction there.

The positive mixed measure of the n floors is bounded above by n\^{}n/n!
times the top measure of their average, by its positive multinomial
expansion. The normalized floor energies of the entries and their
average are uniformly bounded. Section133's bounded all-E1 integrability
estimate, applied to u\_\{r,k\} and that average, gives a uniform bound
for integral(-u\_\{r,k\}) against the mixed floor measure. Therefore its
mass outside W\_k is O(1/k), by summing the n deep-level events. The
increasing winning-set limit has full mass V and the full mixed floor
measures converge to it in total variation, by the same complement
estimate as Section1. It consequently equals the polarization limit. The
same domination and intrinsic integrability show that all original
E1(theta) tests are integrable against this mixed operator, with no mass
on any such test's poles.

For arbitrary simultaneous decreasing E1(theta) nets
u\_\{r,i\}-\textgreater u\_r, all subset averages decrease with E1
endpoints. Top continuity proves mixed weak continuity by the same
finite polarization. If v\_i decreases to v in E1(theta), Section4
applies to each subset average and gives

\[v_i M(u_{1,i},\ldots,u_{n,i})
\longrightarrow vM(u_1,\ldots,u_n).\]

Thus the integrable weighted statement retains the whole intrinsic floor
domain, all original continuous tests and the original arbitrary nets.
The assertion here concerns one common nef theta; mixed products across
unrelated intrinsic nonample classes are not silently supplied.

If V=0, positive bounded mixed products in this common class vanish, as
do every floor energy, top and mixed operator and weighted measure. The
genuine cone may be nonconstant; this says nothing about uniqueness
modulo constants or a positive lower-degree theory.

\subsubsection{6. Both outer derivatives at every intrinsic E1
potential}\label{both-outer-derivatives-at-every-intrinsic-e1-potential}

For continuous h, let F\_theta(h)=E\_theta(P\_theta h). Proposal132
reconstructs genuine bounded nef envelope membership and the supporting
inequality, by passing actual ample supporting inequalities through
bounded operator/energy limits:

\[F_\theta(h+t f)-F_\theta(h)
\le t\int f\,M(P_\theta h),\]

and its counterpart obtained by interchanging the two obstacles. This is
valid for every continuous original f, both signs of t and every
continuous h. It does not assume that the bounded nef envelope is
continuous.

Let u be any intrinsic E1(theta) potential, and take continuous
majorants h\_i decreasing pointwise to u on compact X\^{}an. Such a
directed net exists for every upper-bounded usc function; finite minima
make the family directed. The envelopes p\_i=P\_theta h\_i decrease to
u: their genuine endpoint lies between u and inf h\_i=u. Energy and top
continuity from Section2 give E\_theta(p\_i)-\textgreater E\_theta(u)
and M(p\_i)-\textgreater M(u).

For a fixed real t, q\_i=P\_theta(h\_i+t f) decreases to
q\_t=P\_theta(u+t f). Genuine nef decreasing closure and all-point
obstacle inequalities prove equality with this envelope. Its lower
competitor u-\textbar t\textbar\textbar\textbar f\textbar\textbar{} and
its upper bound u+\textbar t\textbar\textbar\textbar f\textbar\textbar{}
show genuine intrinsic E1 membership. Apply the bounded supporting
inequalities to h\_i,h\_i+t f and pass the original majorant net using
energy and full-C0 operator continuity. For t\textgreater0 this gives

\[\int fM(q_t)\le
\frac{E_\theta(q_t)-E_\theta(u)}t
\le\int fM(u).\]

For t\textless0 the two integral bounds reverse after division. No
continuity of the envelopes themselves was used.

Finally M(q\_t)-\textgreater M(u) weakly against every original
continuous test as t-\textgreater0. In fact
\textbar q\_t-u\textbar\textless=\textbar t\textbar\textbar\textbar f\textbar\textbar{}
wherever u is finite, and both have the same infinite locus. At every
fixed k, the literal floors obey the global finite bound

\[\|\max(q_t,-k)-\max(u,-k)\|_\infty\le |t|\|f\|_\infty.\]

Proposal132's bounded uniform operator continuity proves convergence of
their bounded floor measures. The suprema and normalized negative raw
energies of q\_t are bounded uniformly for small t, by the two order
bounds. Section1's total-variation floor estimate is therefore uniform
in t. First choose k, then let t-\textgreater0. This proves the asserted
full-C0 weak continuity. The two supporting inequalities now squeeze
both signs to

\[\left.\frac d{dt}\right|_{t=0}E_\theta(P_\theta(u+t f))
=\int f\,M(u).\]

This is also the actual outer energy derivative: for an upper-bounded
usc obstacle admitting a genuine intrinsic E1 minorant, the envelopes of
decreasing continuous majorants decrease to its genuine envelope, and
Section2 gives energy convergence. Thus the infimum defining the outer
energy equals the envelope energy used above. At V=0 both sides vanish.

\subsubsection{7. Source applications and remaining original
obligations}\label{source-applications-and-remaining-original-obligations}

The whole final own133 and the whole native121 proof were read before
this continuation. Exact input grants remain confined to the genuine
actual-ample G0 body. The bounded facts are the unreviewed132
construction under BJv4 Theorems7.14/7.18, Lemma7.15, Proposition7.19
and complete nef envelope Theorem5.20; Lemma5.17 is only used on actual
ample classes. Borel locality uses complete actual ample Theorem7.40 and
Lemma7.43, transferred by an explicit finite coefficient argument.

The complete primary Lemma7.22 at literal LF4082--4104 is used only
after unweighted weak convergence and convergence of the integrals of a
nonnegative lsc weight have separately been proved. Proposition7.27 at
LF4175--4208 and complete Lemma7.30/Corollary7.31 at LF4247--4294 are
applied first to bounded actual ample inputs, then to the verified nef
coefficients and finally through proved intrinsic weighted floor limits.
The dimension-only constant convention was read at LF473. Primary URL:
https://arxiv.org/pdf/1801.08229v4 . Whole-book and recursive-foundation
verification are not claimed.

The current result is an unreviewed entire intrinsic nef
operator, mixed product, weighted convergence and differentiability theorem. It does not prove
divisorial weak energy upper semicontinuity and target-integral
continuity on energy sublevels, coercivity and actual variational
attainment, domination or uniqueness for every nef big class. It
constructs no general big nonnef minimal/model reference and no general
nonbig positive lower-degree carrier. Original P16.5 remains open; no
solved-original or independently accepted scoped count changes.

\subsection{IV. Weak energy, intrinsic targets and
coercivity}\label{iv.-weak-energy-intrinsic-targets-and-coercivity}

This is a new unreviewed own continuation of132--134, not an independent
acceptance or complete original P16.5 solution. Active native122 is
frozen and receives none of this material. The actual operator and
intrinsic energy domain below are the original floor constructions; no
auxiliary ample E1 membership of unbounded endpoints is assumed.

\subsubsection{1. Retain the whole intrinsic
domain}\label{retain-the-whole-intrinsic-domain}

Let k,X,L,theta,n and V=theta\^{}n be as in134: arbitrary-cardinality
algebraically closed trivially valued characteristic-zero field, smooth
integral projective X of positive dimension, and any actual nef rational
L, canonical reference0. For the positive-mass assertions assume
V\textgreater0. Every genuine potential is normalized by sup
u=u(v\_triv)=0. Write C(u)=-E\_theta(u)\textgreater=0. All original
E1(theta) functions, including those outside E\textasciitilde1, are
retained.

Use the verified-domain bounded metric pairing and its nef monotonicity
from the primary Theorems7.14/7.18 and Lemma7.15. Genuine normalized nef
compactness and the envelope property come from complete primary
Theorem5.20 and the genuine-cone constructions132/native121. Actual
ample G0 is used solely for actual ample auxiliary classes and their
energy upper semicontinuity.

\subsubsection{2. Divisorial weak upper semicontinuity of the intrinsic
energy}\label{divisorial-weak-upper-semicontinuity-of-the-intrinsic-energy}

Fix an actual ample H, h=c1(H), and an integer k. On genuine theta-psh
scalars with -k\textless=b\textless=0, the published bounded height
pairing gives

\[E_{\theta+\epsilon h}(b)
=\frac1{n+1}\sum_{l=0}^{n+1}\binom{n+1}l\epsilon^l
(\theta,b)^{n+1-l}(h,0)^l.\]

The all-zero reference height is zero. For 1\textless=l\textless=n,
monotonicity replacing each of the m=n+1-l theta-slots between the
constant metrics -k and0 gives

\[-mk\,\theta^{n-l}h^l
\le(\theta,b)^m(h,0)^l\le0.\]

The l=n+1 term is zero. Consequently

\[|E_{\theta+\epsilon h}(b)-E_\theta(b)|
\le\frac{k}{n+1}\sum_{l=1}^n
\binom{n+1}l(n+1-l)\epsilon^l\theta^{n-l}h^l.\]

This tends to zero uniformly over the original bounded domain as
rational epsilon decreases to zero. Each energy on the left auxiliary
actual ample class is upper semicontinuous in divisorial topology by G0;
the inclusion of the genuine theta-cone into that ample cone is
continuous in the same coordinates. A uniform limit of upper
semicontinuous real functions is upper semicontinuous. Thus bounded
E\_theta is weakly upper semicontinuous there.

The genuine floor map u-\textgreater max(u,-k) is continuous in
divisorial topology because it is coordinatewise maximum with a
constant, and genuine endpoints are all-point determined. On the
normalized cone its values lie between -k and0. Hence

\[E_\theta(u)=\inf_{k\ge1}E_\theta(\max(u,-k))\]

is an infimum of upper semicontinuous functions, so is upper
semicontinuous with extended values on the entire normalized genuine
cone. Therefore \{u:sup u=0,C(u)\textless=B\} is a closed subset of the
genuinely compact normalized nef cone, and is weakly compact for every
finite B. This is compactness of actual original functions, not a
completion or a countability assertion.

\subsubsection{3. A quadratic contraction estimate from the actual
metric
products}\label{a-quadratic-contraction-estimate-from-the-actual-metric-products}

For normalized bounded u, set

\[A_l=\int(-u)(\theta+dd^cu)^l\wedge\theta^{n-l},\quad0\le l\le n.\]

The measures are actual positive bounded mixed products. Since
theta\^{}n at the canonical zero metric is V delta\_vtriv and
u(v\_triv)=0, A\_0=0. The actual energy identity says sum
A\_l=(n+1)C(u).

For 0\textless a\textless=1, the genuine convex combination au has
curvature (1-a)theta+a(theta+dd\^{}cu). Its actual mixed expansion gives

\[C(au)=\frac1{n+1}\sum_{r=0}^n\sum_{l=1}^r
\binom rl a^{l+1}(1-a)^{r-l}A_l.\]

Every term is nonnegative. Using a\textsuperscript{\{l+1\}\textless=a}2
and sum\_\{r=l\}\^{}n binom(r,l)=binom(n+1,l+1)\textless=2\^{}\{n+1\},
obtain

\[C(au)\le K_n a^2 C(u),\qquad K_n=2^{n+1}.\]

For every original intrinsic E1 u, apply this to its actual normalized
floors. Their scaled metrics decrease to au and their thresholds aj are
cofinal real reference thresholds. Intrinsic decreasing energy
continuity from134 therefore gives the same bound on the entire floor
domain. This is not an auxiliary-ample estimate on an unverified
unbounded domain.

\subsubsection{4. Exact equivalence of the intrinsic dual finiteness
condition and
integrability}\label{exact-equivalence-of-the-intrinsic-dual-finiteness-condition-and-integrability}

Let mu be a positive original-carrier Radon measure of raw mass V.
Define

\[D_\theta(\mu)=\sup_{u\in E^1(\theta)}
\left(E_\theta(u)-\int u\,d\mu\right).\]

Raw mass V makes this objective invariant under constants. If D is
finite, every genuine intrinsic E1 potential is integrable: after
normalizing u\textless=0 it is bounded above, and integral(-u)=infinity
would make the objective infinite.

Conversely suppose every intrinsic E1(theta) potential is mu-integrable.
The integrals of -u are uniformly bounded on the normalized set
C(u)\textless=1. Otherwise choose u\_j in that set with
integral(-u\_j)\textgreater=2\^{}j. Let s\_m=sum\_\{j=1\}\^{}m
2\^{}\{-j\}u\_j, retaining the remaining convex weight on the genuine0
potential. These actual convex combinations decrease, have generic
value0 and raw energy at least -1 by concavity. Genuine decreasing
closure gives their actual endpoint s, and energy continuity makes s
intrinsic E1. Ordinary countable monotone convergence gives
integral(-s)\textgreater=sum\_j1=infinity, a contradiction. Denote the
finite uniform bound by B\_mu.

For arbitrary normalized u with C=C(u)\textgreater1, choose a=(K\_n
C)\^{}\{-1/2\}. Section3 gives C(au)\textless=1, whence

\[\int(-u)\,d\mu\le B_\mu\sqrt{K_n C}.\]

For C\textless=1 the same integral is at most B\_mu. Thus the objective
-C+integral(-u) is uniformly bounded above, proving
D\_theta(mu)\textless infinity. Its exact intrinsic finiteness condition
is therefore equivalent to integrability of every original intrinsic E1
test. This statement characterizes the condition, not yet the complete
image of the operator.

Every actual top and common-class mixed output constructed in133/134
integrates every such test and has raw mass V. Hence its dual energy is
finite. This direction follows from the already constructed metric
operator, not from an assumed inverse equation.

\subsubsection{5. Coercivity and the remaining attainment
issue}\label{coercivity-and-the-remaining-attainment-issue}

If D=D\_theta(mu)\textless infinity, choose once for all
0\textless a\textless=1/(2K\_n). Applying the dual inequality to the
genuine au and the quadratic estimate gives

\[a\int(-u)\,d\mu\le D+C(au)
\le D+K_n a^2 C(u).\]

Therefore the normalized objective satisfies

\[E_\theta(u)-\int u\,d\mu
\le-\frac12 C(u)+D/a.\]

Every maximizing sequence or net with objective tending to finite D has
uniformly bounded original C-values on a tail. Section2 gives an actual
weakly convergent subnet with intrinsic E1 endpoint. The same inequality
and the integrability argument preserve the exact raw original target
space.

However, upper semicontinuity of the energy alone does not prove upper
semicontinuity of this objective: target integration must be shown
continuous, or a separate strong-Cauchy construction must be completed.
The present argument does not supply that missing passage, actual
maximizer attainment, or uniqueness. General nef big variational
equations therefore remain open here.

\subsubsection{6. Precise source use and next mathematical
gate}\label{precise-source-use-and-next-mathematical-gate}

The bounded height formula, monotonicity and actual mixed operators use
the complete BJv4 Theorems7.14/7.18, Lemma7.15 and Proposition7.19 at
the verified bounded domain. The genuine nef envelope hypothesis is
complete Theorem5.20, not a relabelled ample theorem. The exact G0 body
grants auxiliary actual ample energy usc; the original nef weak usc is
the own uniform coefficient argument above. The whole final own133/134
and native121 Section5 supply only the actual intrinsic floor products,
positivity, concavity and decreasing energy passage reconstructed in
those drafts.

The convex-series idea is compared with the complete printed BJv4
Theorem9.20 beginning at literal LF5266, but the argument above is
independently given on the larger intrinsic nef domain, with a different
explicit quadratic scaling proof. No general nonample result is quoted
from that ample theorem. Primary URL: https://arxiv.org/pdf/1801.08229v4
. No whole-book or recursive foundational proof verification is claimed.

The useful next problem is target-integral continuity on the original
compact energy sublevels, or a complete intrinsic I-Cauchy maximizer
construction, followed by all-point uniqueness for nef big L. Arbitrary
theta-psh test differences need not span all original continuous
functions. Any use of the ample source's PL-theta generation or
divisorial finite-energy Dirac tests must reconstruct their nonample
analogues. Replacing the carrier or assuming unbounded auxiliary-ample
E1 membership would not resolve that gate. General big nonnef references
and general nonbig nonzero numerical data are still not constructed.
Original P16.5 remains open and every assertion in this note requires
independent review.

\subsection{V. Complete independent reconstruction and final
scope}\label{v.-complete-independent-reconstruction-and-final-scope}

PASS for proposals132--135 in the stated nef domain, relative solely to
the inline actual-ample G0 grant. This covers arbitrary-cardinality k,
smooth integral projective X, n\ensuremath{\geq}1, every actual nef rational L and
canonical reference0, and the entire intrinsic floor E1(\ensuremath{\theta}). It does not
accept unchanged original P16.5, an inverse theorem, attainment or
uniqueness. The urn passage documents the ample antecedent used in
independently reconstructing the nef conclusion associated with
premise6; it grants no nef assertion.

Genuine membership was reconstructed as PSH(\ensuremath{\theta})=\ensuremath{\cap}\ensuremath{\epsilon}\textgreater0
PSH(\ensuremath{\theta}+\ensuremath{\epsilon}h). Forward section embeddings preserve the original scalar
because generating canonical section norms have maximum0. For the
reverse inclusion, shifted finite continuous ample FS approximants form
one directed decreasing net: strict comparison with a finite continuous
old minimum gives an open cover, and compactness gives a common later
term. This handles poles and avoids the defective general
extended-valued flattening lemma. Closed intersections inside a
normalized ample cone give genuine compactness and the envelope
property. Lemma5.17 remains ample-only. No unbounded intrinsic endpoint
is placed in E\textasciitilde1; Remark7.16's distinction is retained.

Bounded inputs belong to a common actual ample E1 space. Actual tensor
identities L=(L+H)\ensuremath{-}H construct signed metric products before limits,
with canonical references and continuous dpsh compatibility. At model
stages the coefficient is b\_F\ensuremath{\prod}r(\ensuremath{\theta}r,\ensuremath{\mathcal{X}}+D\_r)\textbar F, with
b\_F=ord\_F(t) and v\_F=b\_F\^{}\ensuremath{-}1 ord\_F\textbar k(X). The collar
identity retains its fixed reference term. Multilinearity proves
auxiliary independence; positive actual ample perturbations give nef
positivity and raw mass V=\ensuremath{\theta}\^{}n.~Borel strict-winning locality follows
by restricting a finite measure-valued polynomial and extracting its
exact coefficients. It does not pass weak convergence through
indicators.

For normalized u, C\_u=\ensuremath{-}E\ensuremath{\theta}(u), positivity of the bounded mixed energy
formula gives \ensuremath{\int}(\ensuremath{-}u\_k)M(u\_k)\ensuremath{\leq}(n+1)C\_u. The increasing strict-winning
restrictions therefore have full mass V and yield
\textbar\textbar M(u\_k)\ensuremath{-}M(u)\textbar\textbar TV\ensuremath{\leq}2(n+1)C\_u/k. Their
limit is Radon and agrees with the bounded operator. Real cofinal floors
and constant shifts preserve it. Expanding the bounded nef height gives
\ensuremath{\int}(\ensuremath{-}v)M(u)\ensuremath{\leq}2\^{}n(C\_u+C\_v), first by bounded-test TV convergence, then
countable monotone convergence. Thus every intrinsic E1 test is
integrable, including endpoints outside E\textasciitilde1.

The normalized I/J and Hölder inequalities are applied first at positive
ample \ensuremath{\epsilon} to bounded actual inputs. Polynomial pairing limits preserve
them with dimension-only constants as V\ensuremath{\epsilon}\ensuremath{\to}V\textgreater0. Ordered floors
satisfy I(a,b)\ensuremath{\leq}n(n+1)(E\ensuremath{\theta}(a)\ensuremath{-}E\ensuremath{\theta}(b))/V. The relative-J estimate gives
uniform vanishing weight tails; bounded-weight TV convergence then
proves fixed and simultaneous intrinsic weighted floor limits. Only
after these limits are established are the inequalities transferred to
intrinsic E1.

For an original decreasing net with E1 endpoint, floor energy continuity
proves full decreasing energy continuity. A common upper shift bounds
all tail costs, making the TV floor error uniform over the original
directed set. Choosing the floor first and the original index second
proves weak convergence on the full original C(X\^{}an). For bounded
discontinuous floors, integral convergence is separately obtained from
the transferred inequality before applying Lemma7.22. Uniform weight
tails handle intrinsic weights and simultaneous decreasing weights. More
generally, for a genuine decreasing psh weight integrable at the
endpoint, polarized intrinsic heights give the lower integral bound;
fixed floors give the upper bound. The resulting integral convergence
establishes Lemma7.22's antecedent before weighted weak convergence. No
arbitrary-Borel-net monotone convergence or restricted test carrier is
used.

Finite subset-average polarization identifies common-\ensuremath{\theta} mixed products as
limits of actual positive bounded mixed floors. Mixed coefficient
locality fixes their restrictions on W\_k=\{all entries\textgreater\ensuremath{-}k\}.
Positive multinomial domination by the average's top product bounds the
complementary mass by O(1/k). Thus their increasing winning-set
construction and polarization limit coincide, with full mass, TV floor
convergence, every intrinsic test integrable, and arbitrary decreasing
mixed weighted-net convergence.

Ample supporting inequalities pass through the bounded nef limits.
Decreasing continuous majorants of any intrinsic E1 base then give
supporting inequalities at that base. The perturbed envelope lies
between u\ensuremath{\pm}\textbar t\textbar\textbar\textbar f\textbar\textbar; its
fixed floors converge uniformly and its floor TV errors are uniform for
small t. Both quotient signs are consequently squeezed to \ensuremath{\int}f M(u). The
same majorant argument identifies envelope energy with outer energy.

The bounded ample-height coefficient estimate is uniform on \ensuremath{-}k\ensuremath{\leq}b\ensuremath{\leq}0.
Hence bounded nef energy is weakly usc; continuous divisorial floor maps
and an infimum establish intrinsic usc. The stated normalized cost sets
\{\ensuremath{-}E\ensuremath{\theta}\ensuremath{\leq}B\} are closed in the genuine compact cone. Positive mixed
expansion, with A0=0 at the canonical generic trivial valuation, proves
C(au)\ensuremath{\leq}2\^{}(n+1)a\ensuremath{^{2}}C(u); cofinal scaled floors extend it intrinsically.
Actual decreasing convex series prove that integrability of all
intrinsic tests bounds their integrals uniformly on C\ensuremath{\leq}1. Scaling then
proves dual finiteness, and every top/common-class mixed output has
finite dual energy. For raw mass V\textgreater0, the objective obeys
E\ensuremath{\theta}(u)\ensuremath{-}\ensuremath{\int}u d\ensuremath{\mu}\ensuremath{\leq}\ensuremath{-}C(u)/2+2\^{}(n+2)D\ensuremath{\theta}(\ensuremath{\mu}). This proves coercivity without
asserting objective usc or attainment. At V=0 all positive common-class
products, floor energies and weighted measures vanish; no probability
normalization or uniqueness is inferred.

Fresh complete pertinent primary units were checked: BJ v4, March2022,
printed pp21,25--26,36--38,42--47,57--63,79--87,91--92; BE v3,
January2021, pp50,55--58,65--66; Fujino--Miyamoto v1 Theorem1.4 and its
pp4--5 proof; Stacks01Q3's complete proposition/proof. Relevant
published BJ units were matched independently. Harmless corrections are
the reversed pairing sign in 7.21(ii), the 7.40 limiting reference to
7.21(ii), and algebraic indices/degrees in 3.31/3.33/7.21/7.30; the n=1
Hölder case uses direct Cauchy--Schwarz. The abstract flattening failure
is substantial, but the supplied finite-continuous nef argument
independently replaces that application. No historical verdict, hash,
nonample note or other campaign supplies a premise.

\end{document}
